\section{Energy modeling}
Contiki-NG estimates energy consumption using a
software-based online mechanism proposed by
\cite{dunkelsSoftwarebasedOnlineEnergy2007}. It uses a linear
model in which each hardware component such as the CPU, the radio transmitter
and the radio receiver, draws a roughly constant current while it is active.

The total energy is then the sum of each component's active time
multiplied by its average current draw and the supply voltage:
\begin{equation}
    E = V \sum_{i} I_i \, t_i ,
\end{equation}
where $V$ is the supply voltage, $I_i$ is the average current draw of
component $i$, and $t_i$ is the time it spent active. Because the
estimation runs entirely in software, it needs no additional
measurement hardware and can easily be ported to different platforms.

Contiki-NG implements this model in its \emph{Energest} module, which
records how long each
component stays in each power state, measured in timer ticks.
Energest tracks give states. Three describe the MCU: active (CPU),
low power mode (LPM),
deep low power mode (DLPM). The other two describe the radio:
transmitting (TX) and receiving (RX). So the energy estimate in
Contiki-NG is calculated as follow. The time $t_x$ spent in each state is
obtained by dividing its tick count by the number of ticks per second
$t_x = \mathit{ticks}_x / \mathit{ENERGEST\_SECOND}$
And the energy consumed by a node is estimated as
\begin{equation}
    E = V \left( I_{\mathrm{CPU}}\, t_{\mathrm{CPU}}
        + I_{\mathrm{LPM}}\, t_{\mathrm{LPM}}
        + I_{\mathrm{DLPM}}\, t_{\mathrm{DLPM}}
        + I_{\mathrm{TX}}\, t_{\mathrm{TX}}
    + I_{\mathrm{RX}}\, t_{\mathrm{RX}} \right),
\end{equation}
where $V$ is the supply voltage, and $I_x$ and $t_x$ are the
average current
draw and the time spent in each state. In many cases where the energy
estimation is used to compare between
homogeneous nodes within the same network, V does not need to be
explicitly calculated to simplify the calculation.

The value of current drawn of each state is gathered from the
datasheet of the MCU and transceiver. Table~\ref{tab:current-draw}
shows the current drawn
of two platforms Z1 \cite{zolertia_z1} and Tmote Sky
\cite{moteiv_tmotesky} with transceiver CC2420 \cite{ti_cc2420}. The
current draw from datasheet should only be taken as a reference
point, which can vary a lot in different situation.

\begin{table}[htbp]
    \centering
    \caption{Average current draw per Energest state for the Zolertia
    Z1 and Tmote Sky platforms.}
    \label{tab:current-draw}
    \begin{tabular}{llcc}
        \hline
        \textbf{Component} & \textbf{State} & \textbf{Z1} &
        \textbf{Tmote Sky} \\
        \hline
        MCU   & Active (CPU)                 & 10\,mA       & 1.8\,mA \\
        & Low-power mode (LPM)         & 23\,$\mu$A   & 54\,$\mu$A   \\
        & Deep low-power mode (DLPM)   & 0\,$\mu$A    & 0\,$\mu$A    \\
        \hline
        Radio & Receive (RX)                 & 18.8\,mA     & 18.8\,mA     \\
        & Transmit (TX)                & 17.4\,mA     & 17.4\,mA     \\
        \hline
    \end{tabular}
\end{table}

% TODO: should I include this
% Time measurement is implemented using the on-chip timers
% of the MSP430. Since the on-chip timers work even when
% the microcontroller is in low-power mode, the time measure-
% ment is non-intrusive. We use the 32768 Hz clock divided
% by 8, producing 4096 clock ticks per second

From the table, radio power is much more significant than MCU power.
Since the removal of duty cycling layer in Contiki-NG in favor of
TSCH, the energy consumption in CSMA mode is dominated by radio
turning for the whole period.
However, TSCH is a much more complicated and memory heavy algorithm
for such old devices. Hence we opt for CSMA and use active CPU ticks
as intermediate value to measure the energy consumption relative to
other nodes in the network.
